Between-donor gene co-expression in single-cell cohorts is usually tested by correlating pseudobulk donor means. We show that this test fails increasingly as cohorts grow: within-donor cell-state covariance leaks into donor means in proportion to each donor's inverse cell count, and log or variance transformations add a factor shared by all genes; each bias is fixed per donor, so their sum grows like the number of donors $D$ while noise grows like $\sqrt D$. In real-cell nulls from 744--1{,}164 PsychAD brain myeloid donors per subtype, pseudobulk called \RfPBmean{} of about 89{,}000 gene pairs per run, and adjusting for cell number and depth, weighting donors or removing principal components left false calls in every run. We propose XDS-M, which measures the two genes of a pair on disjoint random halves of each donor's cells, standardizes by the donor's design only, and averages a studentized cross-product test over five splits; we prove per-pair asymptotic normality. From 100 donors up, XDS-M made a false call in 3--8\% of null runs and met pre-specified bounds in \XDSMholdsLarge{} brain and blood settings; it failed a brain-bank-stratified null in two of three subtypes and passed once bank means were removed. At 25--50 donors only permutation-based XDS was not rejected, with little power. Donor maps show null pseudo-donors ordered by cell count, and a function-text gene map built in the Mantis platform shows pseudobulk's false edges as concentrated within function clusters as real edges, so such coherence does not validate edges. Because Alzheimer's cases and controls differ in cell numbers, pseudobulk also reports false differential co-expression, which a group-centred, unstandardized XDS-M avoids.
